% Appendix source for the grouped Boolean Doob manuscript. Uses amsmath and hyperref.
\section{Concentration required for failure at the stationary endpoint}
\label{sec:endpoint-concentration}

Let $\mu$ be uniform on the binary $L$-cube. Write $P=P_{1/2}$,
$Q=P_{1/\sqrt2}$, $g=Ph$, $\pi=hg$, and $U=h/g$, with $h>0$ and
$E_\mu\pi=1$. Let $K(x,y)=P(x,y)h(y)/g(x)$ and let $\Delta$ denote the
maximum absolute difference of the ordered eigenvalues of $K$ and $P$,
including multiplicities. Assume $\pi\ge b>0$, $\|\pi\|_p\le B$, and exact
uniformity of all stationary coordinate marginals through order $k\to\infty$.
Throughout this section
\[
 p=(17+4\sqrt{13})/9,\qquad q=6p/(p-3),\qquad 1/3=1/p+2/q .
\]
All norms use $\mu$. All constants $b,B$ are uniform along the sequence.

\paragraph{Endpoint concentration criterion.}
Fix rates, independently of dimension, satisfying
\[
 0<\tau<{6\log 2\over q-12},\qquad
 0<\sigma<{6\log 2+12\tau\over p}.
\] Define
\[
 A_k=\{g<e^{-\tau k},\ \pi>e^{\sigma k}\},\qquad
 d_*=\limsup_{k\to\infty}\|\pi{\bf1}_{A_k}\|_p .
\]
If $d_*<b$, then
\[
 \limsup_{k\to\infty}\Delta\ \le\ {B^2 d_*\over (b-d_*)^2}.
 \tag{EC1}
\]
In particular, vanishing of the critical stationary moment on $A_k$
suffices for full spectral recovery at any diverging matching order.
A condition involving the stationary density alone is
\[
 E_\mu[\pi^p{\bf1}_{\{\pi>e^{\sigma k}\}}]\longrightarrow0 .
 \tag{EC2}
\]
Uniform integrability of $\pi^p$ implies (EC2). The converse fails even
among exactly matched chains with a fixed stationary floor.

Since $\tau$ may approach its upper bound, (EC2) suffices for any fixed
$\sigma<6\log 2/(6-p)$ with $\sigma>0$. Indeed
$6q/[p(q-12)]=6/(6-p)$. In particular it is enough that
$\limsup k^{-1}\log_2\|\pi\|_\infty<6/(6-p)$.
For example, when $\pi$ is constant on cosets of a linear map of rank $s$,
each atom of the quotient has probability $2^{-s}$ and
$\|\pi\|_\infty\le B2^{s/p}$. Thus $\limsup s/k<6p/(6-p)$ suffices,
regardless of ambient dimension or the ranks of disjoint column families.

\paragraph{Critical exponents without iteration.}
Put $1/u=1/p+1/q$ and $v=1+4(u-1)$. Substitution of
$9p^2-34p+9=0$ gives $q=4(1+v)-1$.
Forward hypercontractivity for $P$ sends $L^u$ to $L^v$ contractively.
Reverse hypercontractivity with source exponent $-v$ and target exponent
$-q$ gives $\|1/g\|_q\le\|1/h\|_v$: the ratio $(1+v)/(1+q)$ is $1/4$.
The floor implies $1/h=g/\pi\le g/b$. Consequently,
\[
 H:=\|1/g\|_q\le b^{-1}\|g\|_v\le b^{-1}\|h\|_u .
 \tag{EC3}
\]
These inequalities hold in each finite dimension. Simply replacing
$\|h\|_u$ by $BH$ would give no uniform bound.

\paragraph{What matching controls.}
Let $a=E h$, $F=Qh$ and $V=\|F-a\|_2^2=1-a^2$. By Cauchy--Schwarz
under $h\,d\mu$, $1\le a E\pi g$. The high-degree part of $g=QF$
above order $k$ has norm at most $2^{-(k+1)/2}\sqrt V$ in $L^2$.
Matching annihilates the lower nonconstant part against $\pi-1$.
Since $p>2$, it follows that
$V\le a(B+1)2^{-(k+1)/2}\sqrt V$. This includes $V=0$ and yields
$\|g-1\|_2\le C_B2^{-k/2}$ using $1-a\le V$.
In fact $\|g-1\|_3\le C_B2^{-k/2}$ because $g-a=Q(F-a)$ and $Q$
is contractive from $L^2$ to $L^3$.

\paragraph{A low denominator forces high denominators among its neighbors.}
The floor gives $g=P(\pi/g)\ge bP(1/g)$. For $T\ge T_0=4/b$, put
$S_T=\{g>bT/2\}$. If $g(x)^{-1}>T$, then
\[
 P_x(S_T^c)\le(bT/2)P(1/g)(x)<1/2 .
\]
The cubic estimate for $g-1$ gives
$\mu(S_T)\le C_{B,b}2^{-3k/2}T^{-3}$. For any $1<r<4$, choose
$m=3/(4-r)>1$ and $\ell=4m-3$. Forward hypercontractivity of $P$
from $L^m$ to $L^\ell$ and Markov's inequality give
\[
 \mu\{g^{-1}>T\}\le 2^\ell\|P{\bf1}_{S_T}\|_\ell^\ell
 \le 2^\ell\mu(S_T)^r
 \le C_{B,b,r}\,2^{-3rk/2}T^{-3r}.                 \tag{EC4a}
\]
This additional use of the floor controls the geometry of the exceptional
set before applying (EC3).

Write $D=\{g<e^{-\tau k}\}$. Choose $r<4$ close enough to four that both
$(q-3r)\tau<(3r/2)\log 2$ and
$p\sigma<(3r/2)\log 2+3r\tau$. This is possible by the two strict rate
conditions and $q>12$. Integrating (EC4a) up to $T=e^{\tau k}$ bounds
the integral of $g^{-q}$ on $\{T_0<1/g\le e^{\tau k}\}$ by
$C2^{-3rk/2}(1+e^{(q-3r)\tau k})=o(1)$.
On $\{1/g\le T_0\}$ the reciprocal is bounded and converges to one in
probability. Consequently
\[
 \|g^{-1}{\bf1}_{D^c}\|_q\to1,\qquad
 \|(g^{-2}-1){\bf1}_{D^c}\|_{q/2}\to0 .           \tag{EC4}
\]
On $D\setminus A_k$, $\|\pi{\bf1}_{D\setminus A_k}\|_p
 \le e^{\sigma k}\mu(D)^{1/p}=o(1)$ by (EC4a) at $T=e^{\tau k}$.
Hence $d_*=\limsup\|\pi{\bf1}_D\|_p$.

\paragraph{Absorption on the exceptional set.}
Split $h=\pi/g$ over $D^c$ and $D$ in (EC3). H\"older gives
\[
 bH\le B\|g^{-1}{\bf1}_{D^c}\|_q+\|\pi{\bf1}_D\|_p H .
\]
If $d_*<b$, then $\limsup H\le B/(b-d_*)$ by (EC4). A second application
of H\"older gives $\|U{\bf1}_D\|_3\le\|\pi{\bf1}_D\|_p H^2$.
Also $\|\pi{\bf1}_D\|_3\le B\mu(D)^{1/3-1/p}\to0$ and, by (EC4),
$\|(U-\pi){\bf1}_{D^c}\|_3\to0$. This proves the bound on the right of
(EC1) for $\limsup\|U-\pi\|_3$.

For completeness, $\|QM_{U-\pi}Q\|\le\|U-\pi\|_3$ by forward
$2$-to-$3$ hypercontractivity. For $M_{\pi-1}$ split each factor of $Q$
at Walsh degree $j=\lfloor k/2\rfloor$. Matching kills the low-low term.
Writing $p'=p/(p-1)$, the remaining high factors have $2$-to-$2p'$ norm
at most $\theta^{j+1}$ where $\theta=\sqrt{(2p'-1)/2}<1$: factor
$Q=P_{1/\sqrt{2p'-1}}P_\theta$ and apply the Walsh tail bound to $P_\theta$.
The same factorization bounds the low factors by one. The three terms
therefore contribute at most $3(B+1)\theta^{j+1}$. Finally $QM_UQ$ and
the symmetric conjugate of $K$ are the two Gram products of
$M_{\sqrt U}Q$. Min--max gives (EC1).

\paragraph{Necessary joint concentration.}
If $\Delta\ge\eta>0$ along a subsequence, then for \emph{every} fixed
admissible pair $\tau,\sigma$ along that subsequence,
\[
 \liminf\|\pi{\bf1}_{A_k}\|_p\ge t_\eta,\qquad
 B^2t_\eta=\eta(b-t_\eta)^2,\quad 0<t_\eta<b .
\]
Indeed the right-hand side of (EC1) is strictly increasing from zero to
infinity on $[0,b)$. A further subsequence below $t_\eta$ would contradict
(EC1). This localizes both the large stationary density and the small
denominator to the same vertices. The bound is on the integral of the
$p$th power; their stationary probability can tend to zero.

\paragraph{A separating example.}
Take $n=\lfloor\sqrt{k}\rfloor$ disjoint blocks of $k+1$ bits. Let $S$
be the event that every block parity is positive, $m=2^{-n}$, and prescribe
\[
 \pi=1+m^{-1/p}({\bf1}_S-m).
\]
This density has mean one and a fixed positive floor for all sufficiently
large $k$. Its nonconstant Walsh terms have degree at least $k+1$ by the
block-parity expansion. It is realizable as $hPh$ by positive symmetric
diagonal scaling of the strictly positive matrix $P$ with row sums $\pi(x)$.
Here $\log\|\pi\|_\infty=O(\sqrt{k})$, so (EC2) holds. Yet
$E[\pi^p{\bf1}_S]\to1$ while $\mu(S)\to0$, and
$E[\pi^{p+\epsilon}{\bf1}_S]\to\infty$ for every fixed $\epsilon>0$.
The complement contributes a bounded critical moment. This example
separates the hypotheses on stationary integrability; the theorem allows
general stationary laws beyond this particular block geometry.

\paragraph{Relation to earlier work and review context.}
The forward and reverse inequalities are classical inputs
(\href{https://doi.org/10.5802/aif.357}{Bonami};
\href{https://arxiv.org/abs/1108.1210}{Mossel, Oleszkiewicz and Sen}).
The realization in the example uses
\href{https://www.cambridge.org/core/services/aop-cambridge-core/content/view/CDCA18916EAC355292A931C731AAC86C/S0008414X00040293a.pdf}{Sinkhorn's prescribed-row-sum theorem}.
The programme's
\href{https://airr.science/papers/ARR-2026-4T35VQ6MZB8F99F4/versions/v3/}{stationary-moment paper}
and its subsequent matching-order estimate supply the setting and the
matching-energy argument. The new step is absorption at the exact fixed
exponent using the critical mass on the common exceptional set. It gives
a conditional endpoint theorem. Whether the original floor, moment bound
and matching alone force recovery is open.


\section{A local return test for the full ordered spectrum}
\label{sec:local-return-test}

Keep the half-noise model, the uniform matching through order $k$, the
stationary floor $b$ and the moment bound $B$ of the preceding section.
In this section any fixed finite $p>3$ is allowed. Put $q=6p/(p-3)$,
choose $0<\sigma<\tau<(\log 2)/q$, and set
\[
 D=\{g<e^{-\tau k}\},\qquad A=D\cap\{\pi>e^{\sigma k}\}.
\]
Write $\nu=\pi\mu$, $s=\|{\bf1}_A K{\bf1}_A\|_{2\to2,\nu}$ and
\[
 R_A=\max_{\varnothing\ne E\subseteq A}
       {\Pr}_{\nu}(X_1\in E\mid X_0\in E).
\]
Both quantities are zero for empty $A$. There exist $c,C>0$, depending
only on the displayed fixed parameters, for which
\[
 s/4-Ce^{-ck}\le\Delta\le s+Ce^{-ck},\qquad
 |\lambda_2(K)-\max(1/2,s)|\le Ce^{-ck},\qquad
 R_A\le s\le\sqrt{2R_A-R_A^2}.                         \tag{LR1}
\]
In particular the full ordered spectrum converges precisely when every
subset of this specified exceptional set has vanishing one-step return,
uniformly over those subsets. The second eigenvalue can converge to $1/2$
even when the criterion for the full spectrum fails; the formula for
$\lambda_2$ by itself only detects $s$ above $1/2$. The two-sided bound
on $\Delta$ does not require knowing which noise eigenvalue is displaced.

There is a more detailed reconstruction with multiplicities. Form a list
from the ordinary cube eigenvalues (each $2^{-j}$ occurs $\binom Lj$ times)
and all eigenvalues of the killed matrix $K_{AA}$ on $A$. Sort in decreasing
order and keep the first $2^L$ entries. Each entry differs from the
corresponding eigenvalue of $K$ by at most $Ce^{-ck}$. The additional
matrix has order $|A|\le C_B2^{L-k}$, with an empty matrix allowed.

\paragraph{Isolating the multiplier that survives matching.}
The matching-energy proof above gives $\|g-1\|_2\le C_B2^{-k/2}$
also for this $p$. The floor and Jensen give $1/g\le Pg/b$, and
$Pg=P_{(1/2)^{3/2}}Qh$ with $\|Qh\|_2=1$. Forward hypercontractivity
therefore gives $\|1/g\|_9\le1/b$.
On $D^c\cap\{g<1/2\}$, $\mu\{g<1/2\}\le C_B2^{-k}$ and
$g^{-2}\le e^{2\tau k}$. Since $q\tau<\log 2$,
the $L^{q/2}$ norm of $(g^{-2}-1)$ on that part decays exponentially.
On $\{g\ge1/2\}$ the function $g^{-2}-1$ is bounded, and its absolute
value is at most $16|g-1|$; interpolation with the $L^2$ estimate gives
exponential decay in $L^{q/2}$ there too. Thus
$\|(U-\pi){\bf1}_{D^c}\|_3\le Ce^{-ck}$ by $1/3=1/p+2/q$.

On the remaining low-denominator vertices outside $A$, use the ninth
inverse moment directly:
\[
 \|U{\bf1}_{D\setminus A}\|_3^3
 \le e^{3\sigma k}\int_D g^{-6}\,d\mu
 \le b^{-9}e^{-3(\tau-\sigma)k}.
\]
Also $\|\pi{\bf1}_D\|_3\le B\mu(D)^{1/3-1/p}\le Ce^{-ck}$.
Together these bounds give $\|U{\bf1}_{A^c}-\pi\|_3\le Ce^{-ck}$.
The signed low/high Walsh estimate for $\pi-1$ above uses only $p>3$.
By $2$-to-$3$ hypercontractivity,
\[
 \|QM_UQ-(P+S)\|\le Ce^{-ck},\qquad S=QM_{U{\bf1}_A}Q .       \tag{LR2}
\]
The two Gram products of $M_{\sqrt U{\bf1}_A}Q$ identify $\|S\|=s$.
In particular $S$ is positive semidefinite with $s\le1$.

\paragraph{Separation from the ordinary low eigenmodes.}
Let $\Pi$ project onto Walsh degrees at most $m=\lfloor\delta k\rfloor$,
where $0<\delta<1$ is fixed with $\delta\log 7<(\log 2)/12$.
The floor bound above gives $\|U\|_{3/2}\le B/b^2$ since
$1/p+2/9<2/3$. For a unit vector $z$ in this low space,
$\|Qz\|_8\le7^{m/2}$: apply the ordinary $2$-to-$8$ inequality
to a polynomial with each degree $j$ multiplied by $7^{j/2}$.
H\"older with reciprocal exponents $2/3,1/4,1/12$ gives
\[
 \langle z,Sz\rangle \le (B/b^2)\,7^m\,\mu(A)^{1/12}
 \le Ce^{-ck}.
\]
Positivity bounds the norm of the off-diagonal block of $S$ by
$\|\Pi S\Pi\|^{1/2}\|S\|^{1/2}$. On the complementary space $P$ has norm
at most $2^{-(m+1)}$. Consequently $QM_UQ$ differs in operator norm by
at most $Ce^{-ck}$ from the block diagonal operator
\[
 P|_{\operatorname{ran}\Pi}\ \oplus\
 S_H,\qquad S_H=(I-\Pi)S(I-\Pi)|_{\operatorname{ran}(I-\Pi)} .
\]
Its high-block norm $s_H$ differs from $s$ by at most $Ce^{-ck}$ as well.
For large $k$, the low block contains $1$ once and $1/2$ with its ordinary
multiplicity. The high block is positive with top eigenvalue $s_H\le1$.
This proves the second-eigenvalue estimate by min--max and Gram conjugation.
The upper bound on $\Delta$ follows directly from (LR2).

To verify the reconstruction, work on two copies of the cube space and
compare $P\oplus S$ to $(\Pi P\Pi)\oplus((I-\Pi)S(I-\Pi))$.
The norm error is exponentially small by the same two block estimates.
After permuting subspaces the latter is the displayed block diagonal
operator with $2^L$ extra zeros. Also $S$ has the nonzero eigenvalues of
$K_{AA}$ with their multiplicities, by the Gram identity. Thus min--max
on the doubled space identifies the first $2^L$ sorted entries with the
claimed precision. The size bound follows from $\mu(A)\le C_B2^{-k}$.

For the lower bound, if $s_H\le4\,2^{-m}$ it is absorbed by the error.
Otherwise choose $j$ with $s_H/4<2^{-j}\le s_H/2$, so $1\le j\le m$.
In the ordered noise spectrum, let $i$ be the first index at level $2^{-j}$.
The low block already supplies the $i-1$ eigenvalues strictly above this
level, all at least $2^{1-j}$. The high block supplies one more at least
$2^{1-j}$. Its addition raises the $i$th eigenvalue by at least $2^{-j}$.
This proves the first bound despite the growing noise multiplicities.
Enlarging $C$ covers the finitely many smaller matching orders.

\paragraph{A test phrased entirely in one-step probabilities.}
Testing with an indicator gives $R_A\le s$. For the reverse bound take a
nonnegative unit top vector $f$ for the killed chain, extended by zero
outside $A$. Under its full stationary pair law $J$, integration over the
level sets of $f^2$ gives
\[
 \tfrac12 E_J|f(X)^2-f(Y)^2|\ge(1-R_A)\|f\|_{2,\nu}^2.
\]
Indeed every nonempty positive level set lies in $A$, and its outward
flow is at least $1-R_A$ times its stationary mass. Cauchy--Schwarz on
$(f(X)-f(Y))(f(X)+f(Y))$ bounds the left side by $\sqrt{1-s^2}$
using stationarity and reversibility. This proves the last part of (LR1).

\paragraph{What nearly complete trapping requires in two ordinary noise steps.}
For $x\in A$, the floor bound $P(1/g)(x)\le g(x)/b$ implies
$P_x(A)\le e^{-2\tau k}/b$. Yet the stationary conditional return under
$K$ may be much larger. For nonempty $E\subseteq A$, Cauchy--Schwarz
under $P_x$ and the floor give
\[
 K_x(E^c)\ge {b\over g(x)}P_x({\bf1}_{E^c}/g)
 \ge {b\,P_x(E^c)^2\over g(x)\,Pg(x)} .
\]
If $\lambda_2(K)\to1$, (LR1) supplies maximizing sets $E_k$ with
conditional return tending to one. Averaging this last bound for $x\in E_k$
gives $E_{\nu(\cdot\mid E_k)}[(g\,Pg)^{-1}]\to0$; hence $g\,Pg$ diverges
in probability under those conditional laws. Conversely a uniform upper
bound $g\,Pg\le M$ on $A$ gives $s\le1-b(1-o(1))/M$ by the maximum
killed row sum, for nonempty $A$. This rules out convergence of the
second eigenvalue to one in that regime.

\paragraph{A family with a correct spectral gap and incorrect lower bands.}
At half noise the
\href{https://github.com/lluiseriksson/THE-ERIKSSON-PROGRAMME/blob/7ac6acd4fe2dcf992c6acc11eb406b8a8323d7c3/papers/stationary-weighting354/research/2026-11-18-small-support/ALL-CORRELATIONS-SHELL-TRAP.md}{programme's matched shell construction}
provides some fixed $p_0>3$, a fixed floor, $k=\lfloor\beta L\rfloor$,
$\|\pi_{\rm trap}-1\|_{p_0}\le e^{-c_0 L}$, and $\lambda_2\to1$.
Fix any $0<a_0<1$. Interpolate the stationary densities by
$\pi_t=1+t(\pi_{\rm trap}-1)$, $0\le t\le1$. They share all the matching,
floor and moment estimates with uniform constants. At each fixed dimension
positive symmetric diagonal scaling gives the unique $h_t>0$ satisfying
$h_t P h_t=\pi_t$ (here $P h_t$ is evaluated first). This scaling is
continuous in $t$: each $h_t(x)$ is bounded above by
$\sqrt{\sup_t\pi_t(x)/P(x,x)}$, and the floor divided by the largest such
upper bound bounds every coordinate below. Any subsequential limit solves
the same scaling equations, so uniqueness gives continuity.

Set $N(t)=\|QM_{h_t/(P h_t)}Q-P\|$. This continuous quantity is zero at
$t=0$. By (LR2), $|N(t)-s(t)|\le Ce^{-ck}$ uniformly in $t$, even though
the exceptional vertex set can change discontinuously. By (LR1),
$s(1)\to1$. The intermediate value theorem gives $t_L$ with $N(t_L)=a_0$
for every sufficiently large $L$. For these chains
\[
 |\lambda_2(K)-\max(1/2,a_0)|\le Ce^{-ck},\qquad
 \Delta\ge a_0/4-Ce^{-ck}.
\]
Moreover for every fixed $j$ with $2^{-j}>a_0$, all the ordered eigenvalues
in degrees $0$ through $j$, counted with the reference multiplicities,
are within $Ce^{-ck}$ of their reference values. Indeed the low block
contains all those eigenvalues exactly, and the entire high block is below
$2^{-j}$ for large $L$. In particular any prescribed finite number of leading
Walsh bands can be asymptotically correct while the full ordered spectrum
has a fixed error. The available $p_0$ comes from the matched shell result
below the stationary endpoint. This interpolation constructs no counterexample
at the endpoint itself.

\paragraph{Scope and provenance.}
The coarea estimate is the Dirichlet form argument underlying
\href{https://arxiv.org/abs/math/0505690}{Goel, Montenegro and Tetali's spectral profile}.
The low-degree norm bound uses classical Boolean hypercontractivity.
The contribution here is the exponential operator approximation on a
specific joint tail and the eigenvalue count that converts a killed mode
there into a defect of the full ordered spectrum. The unrestricted
endpoint still requires controlling returns for all subsets of that tail.
